%% -*- coding: utf-8 -*-
\documentclass[12pt,pagesize,paper=192mm:108mm]{scrbook} 
%1920x1080 1280x720
\areaset[current]{192mm}{108mm}
\usepackage{calc}
\usepackage[T2A]{fontenc}
\usepackage[utf8]{inputenc}
\usepackage[english,russian]{babel}
\usepackage{microtype}
\usepackage{misccorr}
\usepackage{cmap}
%\usepackage[unicode=true]{hyperref}
\usepackage{graphicx}
\usepackage{amssymb}
\usepackage{amsmath}
%\usepackage{srcltx}
\usepackage{textcomp}
\usepackage{xspace}
%научные символы и смайлики \smiley \frownie
\usepackage{wasysym}
\usepackage{ccicons}
\begin{document}
\begin{titlepage}
  \vspace*{-0.5em}
  \begin{center}    
    \hspace*{3em}
    \begin{minipage}[t]{3em}
      \includegraphics[width=\textwidth]{../BINP-logo}
    \end{minipage}\hfill
    \begin{minipage}{0.23\linewidth}
    \includegraphics[width=\textwidth]{../NSU-logo}
    \end{minipage}
    \hfill
    \hspace*{6em}

    Кафедра теоретической физики физического факультета НГУ
    \medskip

    \Large
    Профессор Черняк В.\,Л.
    \bigskip

    \huge
    \textbf{Теория электрослабых взаимодействий}
    \bigskip

    \Large
    Лекция № 9
    \vfill

    % \normalsize
    % \begin{minipage}{0.65\linewidth}
    % \end{minipage}
    \vfill

\normalsize    Новосибирск 2013
  \smallskip

  \ccbysa
  \end{center}
\end{titlepage}
\newpage

\vspace*{-1em}
\begin{center}
 \vfill
  \begin{minipage}{0.66\linewidth}
    Анализ формулы для ширины распада мюона на электрон, электронное
    антинейтрино и мюонное нейтрино. Полная ширина распада
    $\tau$"=лептона.  Анализ формулы для лептонной ширины распада
    $\pi$"=мезона. Распады псевдоскалярных мезонов ($\pi^-$, $K^-$,
    $D^-$, $B^-$) на~лептонную пару: $f_{\pi}$ и $f_K$, отношение
    ширин распада $K^-$ в лептон и антинейтрино и $\pi^-$ в лептон и
    антинейтрино, предел $SU(3)$ изотопической симметрии. Распад
    заряженного векторного мезона на лептонную пару (лептон и
    нейтрино).  Распад $\rho^-$. Константа $f_{\rho}$. Отсутствие
    подавления по спиральности.  Трудности экспериментального
    наблюдения слабого распада $\rho^-$. Распад $\tau$"=лептона
    на~$\pi$/$\rho$ и нейтрино~($\nu_{\tau}$).
  \end{minipage}
  \vfill
  % Новосибирск 2013

  % \normalsize \ccbysa
\end{center}
\end{document}
